%% kapa: exam=1 round=26 year=2015 period=2 form=A subject=appraisal_law subject_ko=부동산관계법규 seq=6 number=6
%% source: Q-Net 감정평가사 기출문제 (cst00309) artlSeq=5138001 file=extracted/1008/5138001/제26회_감정평가사 1차_기출/제26회_감정평가사 1차_기출/2교시_감평_A형.pdf
%% conversion: kordoc 4.24.51 -> tools/gen1.py (auto); review=proofread-sonnet; flags=none
\begin{question}{6}
국토의 계획 및 이용에 관한 법령상 이행강제금의 부과에 관한 설명으로 옳지 않은 것은?
\begin{choices}
  \item 토지의 이용의무를 이행하지 아니한 자에 대하여 이용의무의 이행을 명하는 경우, 그 이행명령은 3월 이내의 기간을 정하여 문서로 하여야 한다.
  \item 시장ㆍ군수 또는 구청장은 이행강제금을 부과하기 전에 이행강제금을 부과ㆍ징수한다는 뜻을 미리 문서로 알려야 한다.
  \item 이행강제금 부과처분을 받은 자가 이의를 제기하려는 경우에는 부과처분의 고지를 받은 날부터 30일 이내에 이의를 제기하여야 한다.
  \item 토지의 이용의무에 대한 이행명령을 받은 자가 그 명령을 이행하는 경우에는 명령을 이행하기 전에 이미 부과된 이행강제금은 징수하여서는 아니된다.
  \item 토지의 이용 의무기간이 지난 후에는 이행강제금을 부과할 수 없다.
\end{choices}
\end{question}
\answer{4}
